% unit: u31
% sources: sources/074-positive-definite-matrices-and-minima.txt, sources/053-lecture-27-positive-definite-matrices-and-minima.txt, sources/212-mit18-06scf11-ses3-3sum-pdf-download-file.txt, sources/054-problem-solving-positive-definite-matrices-and-minima.txt, sources/213-mit18-06scf11-ses3-3prob-pdf-download-file.txt, sources/214-mit18-06scf11-ses3-3sol-pdf-download-file.txt
\section{Positive definite matrices and minima}\label{sec:u31}

Section~\ref{sec:u29} defined a positive definite matrix as a symmetric matrix with positive
eigenvalues and showed that positive pivots, or positive upper-left determinants, are
equivalent tests. This unit adds the quantity that explains \emph{why} these matrices matter:
the quadratic expression $\vb{x}\T A\vb{x}$. It gives a fourth test, which most books take as the
definition, and it connects linear algebra to calculus (a minimum of a function of several
variables) and to geometry (ellipses and ellipsoids). Pivots from the first week,
determinants from the middle of the course and eigenvalues from the last weeks all meet here.
All matrices in this unit are real and symmetric.

\subsection{The energy \texorpdfstring{$\vb{x}\T A\vb{x}$}{xTAx} and four equivalent tests}

\begin{definition}[Quadratic form; positive (semi)definite]\label{def:u31-form}
The \emph{quadratic form} of a symmetric $n\times n$ matrix $A$ is
$f(\vb{x})=\vb{x}\T A\vb{x}=\sum_{i,j}a_{ij}x_ix_j$. The matrix is \emph{positive definite} if
$\vb{x}\T A\vb{x}>0$ for every $\vb{x}\ne\vb{0}$, and \emph{positive semidefinite} if $\vb{x}\T A\vb{x}\ge0$
for every $\vb{x}$.
\end{definition}

For a $2\times2$ matrix the form is short enough to read off:
\[
  \begin{bmatrix}x&y\end{bmatrix}\begin{bmatrix}a&b\\b&c\end{bmatrix}\begin{bmatrix}x\\y\end{bmatrix}
  =ax^2+2bxy+cy^2 .
\]
The diagonal entries multiply the pure squares and each off-diagonal pair produces a cross term.
The form is purely quadratic: no linear terms, no constant. The squares $ax^2$ and $cy^2$ are
positive when $a,c>0$; the cross term $2bxy$ can have either sign. Positive definiteness means
the squares always win.

\begin{theorem}[Four tests]\label{thm:u31-tests}
For a symmetric matrix $A$ the following are equivalent; any one of them makes $A$ positive
definite.
\begin{enumerate}[(1)]
  \item All eigenvalues are positive: $\lambda_i>0$.
  \item All upper-left determinants are positive: $D_1>0,\dots,D_n>0$.
  \item All pivots are positive (elimination without row exchanges): $d_i>0$.
  \item $\vb{x}\T A\vb{x}>0$ for every $\vb{x}\ne\vb{0}$.
\end{enumerate}
For $A=\left[\begin{smallmatrix}a&b\\b&c\end{smallmatrix}\right]$ these read: $\lambda_1,\lambda_2>0$;
$a>0$ and $ac-b^2>0$; pivots $a>0$ and $(ac-b^2)/a>0$; $ax^2+2bxy+cy^2>0$ unless $x=y=0$.
\end{theorem}
\begin{proof}
Tests (1), (2), (3) were shown equivalent in Section~\ref{sec:u29}. It remains to connect (4)
with (1). Write $A=Q\Lambda Q\T$ (spectral theorem) and $\vb{y}=Q\T\vb{x}$. Then
\[
  \vb{x}\T A\vb{x}=\vb{y}\T\Lambda\vb{y}=\lambda_1y_1^2+\lambda_2y_2^2+\cdots+\lambda_ny_n^2 .
\]
If every $\lambda_i>0$, this is positive whenever $\vb{y}\ne\vb{0}$, that is whenever $\vb{x}\ne\vb{0}$
($Q$ is invertible). Conversely, if (4) holds and $A\vb{x}=\lambda\vb{x}$ with $\vb{x}\ne\vb{0}$, then
$0<\vb{x}\T A\vb{x}=\lambda\,\vb{x}\T\vb{x}$, so $\lambda>0$.
\end{proof}

The same computation shows that $A$ is positive \emph{semi}definite exactly when all its
eigenvalues are $\ge0$. A semidefinite matrix that is not definite is singular: it has the
eigenvalue $0$, fewer than $n$ pivots, and a nonzero $\vb{x}$ with $\vb{x}\T A\vb{x}=0$.

\begin{example}[One matrix, three regimes]\label{ex:u31-family}
For which $c$ is $A_c=\begin{bmatrix}2&6\\6&c\end{bmatrix}$ positive definite? We do not want
eigenvalues; the determinant test is fastest. $D_1=2>0$ always, and $D_2=2c-36>0$ exactly when
$c>18$. Its form is $f_c(x,y)=2x^2+12xy+cy^2$.
\begin{itemize}
  \item $c=18$ is the borderline: $\det A_{18}=0$, so $A_{18}$ is singular with eigenvalues
        $0$ and $20$ (the trace). It has one pivot, $2$, because row~2 is $3$ times row~1. Its form
        is $2x^2+12xy+18y^2=2(x+3y)^2\ge0$, which vanishes along the line $x=-3y$, for instance at
        $(3,-1)$. $A_{18}$ is positive semidefinite.
  \item $c=7$ fails badly: $\det A_7=14-36=-22$, the pivots are $2$ and $-22/2=-11$, and the
        eigenvalues (trace $9$, determinant $-22$) are $11$ and $-2$. The form goes negative:
        $f_7(1,-1)=2-12+7=-3$. Its graph is a \emph{saddle}: up along some directions, down along
        others.
  \item $c=20$ passes: $D_2=40-36=4>0$, the pivots are $2$ and $4/2=2$, and the eigenvalues
        have product $4$ and sum $22$, so both are positive. The graph of $f_{20}$ is a bowl.
\end{itemize}
\end{example}

\begin{trap}
Checking $\vb{x}\T A\vb{x}>0$ at a few convenient vectors proves nothing. The indefinite $A_7$ gives
$f_7(1,0)=2$, $f_7(0,1)=7$ and $f_7(1,1)=21$, all positive, yet $f_7(1,-1)=-3$. Positive diagonal
entries are necessary but far from sufficient: the cross term must be outweighed in every
direction, which is what the determinant $ac-b^2>0$ measures.
\end{trap}

\subsection{Completing the square is elimination}

How can we see that $f_{20}=2x^2+12xy+20y^2$ is positive everywhere except the origin? Write it
as a sum of squares. Put all the $x$'s into one square: $2(x+3y)^2=2x^2+12xy+18y^2$ produces the
first two terms correctly and $18y^2$. We have $20y^2$, so two are left over:
\[
  2x^2+12xy+20y^2=2(x+3y)^2+2y^2 ,
\]
a sum of squares with positive weights, zero only when $y=0$ and $x+3y=0$. For $c=7$ the same
step leaves $7-18=-11$:
\[
  2x^2+12xy+7y^2=2(x+3y)^2-11y^2 ,
\]
negative for example at $(x,y)=(-3,1)$, where it equals $0-11=-11$. The number $18$ is exactly the
borderline found by the determinant test.

The numbers $2$, $3$, $2$ are not new. Elimination on $A_{20}$ subtracts $3$ times row~1 from row~2:
\[
  \begin{bmatrix}2&6\\6&20\end{bmatrix}\longrightarrow\begin{bmatrix}2&6\\0&2\end{bmatrix},
  \qquad
  L=\begin{bmatrix}1&0\\3&1\end{bmatrix},\quad D=\begin{bmatrix}2&0\\0&2\end{bmatrix}.
\]
The pivots stand outside the squares and the multiplier $3$ sits inside.

\begin{proposition}[Sum of squares from \texorpdfstring{$A=LDL\T$}{A = LDLT}]\label{prop:u31-ldl}
If $A=LDL\T$ (elimination without row exchanges, pivots $d_1,\dots,d_n$), then
\[
  \vb{x}\T A\vb{x}=d_1\,(\boldsymbol{\ell}_1\T\vb{x})^2+d_2\,(\boldsymbol{\ell}_2\T\vb{x})^2+\cdots+d_n\,(\boldsymbol{\ell}_n\T\vb{x})^2 ,
\]
where $\boldsymbol{\ell}_i$ is column $i$ of $L$. In particular positive pivots give
$\vb{x}\T A\vb{x}>0$ for $\vb{x}\ne\vb{0}$, which proves test (3)$\Rightarrow$(4) directly.
\end{proposition}
\begin{proof}
$\vb{x}\T LDL\T\vb{x}=\vb{z}\T D\vb{z}$ with $\vb{z}=L\T\vb{x}$, and component $i$ of $\vb{z}$ is
$\boldsymbol{\ell}_i\T\vb{x}$. So $\vb{z}\T D\vb{z}=\sum d_iz_i^2$. If all $d_i>0$ this vanishes only for
$\vb{z}=\vb{0}$, and since $L\T$ is invertible (ones on its diagonal) that means $\vb{x}=\vb{0}$.
\end{proof}

\begin{example}[A \texorpdfstring{$3\times3$}{3 by 3} example in every test]\label{ex:u31-K}
Take the second-difference matrix
\[
  K=\begin{bmatrix}2&-1&0\\-1&2&-1\\0&-1&2\end{bmatrix}.
\]
\begin{itemize}
  \item \emph{Determinants:} $D_1=2$, $D_2=4-1=3$, $D_3=2\cdot3-(-1)(-2)=4$. All positive.
  \item \emph{Pivots:} $D_k/D_{k-1}$ gives $2,\ \tfrac32,\ \tfrac43$. Elimination uses multipliers
        $\ell_{21}=-\tfrac12$ and $\ell_{32}=-\tfrac23$.
  \item \emph{Eigenvalues:} $2-\sqrt2,\ 2,\ 2+\sqrt2$, all positive. Checks: their sum is $6=\tr K$
        and their product is $(4-2)\cdot2=4=\det K$.
  \item \emph{Energy:} $\vb{x}\T K\vb{x}=2x_1^2+2x_2^2+2x_3^2-2x_1x_2-2x_2x_3$, and completing squares with
        the pivots outside and the multipliers inside,
        \[
          \vb{x}\T K\vb{x}=2\bigl(x_1-\tfrac12x_2\bigr)^2+\tfrac32\bigl(x_2-\tfrac23x_3\bigr)^2+\tfrac43x_3^2 .
        \]
        Expanding: $2x_1^2-2x_1x_2+\tfrac12x_2^2$, plus $\tfrac32x_2^2-2x_2x_3+\tfrac23x_3^2$, plus
        $\tfrac43x_3^2$, which totals the form above.
\end{itemize}
Every test agrees: $K$ is positive definite, and the graph of $\vb{x}\T K\vb{x}$ over $\R^3$ is a bowl in
four dimensions.
\end{example}

Figure~\ref{fig:u31-levels} shows what separates a bowl from a saddle, by drawing curves
$f(x,y)=k$ on which the form is constant.

\begin{figure}[ht]\centering
\begin{tikzpicture}[scale=0.72,>=Stealth]
  \begin{scope}
    \draw[->,gray] (-3.6,0) -- (3.6,0) node[right] {$x$};
    \draw[->,gray] (0,-1.9) -- (0,1.9) node[above] {$y$};
    \draw[thick,blue,domain=0:360,samples=90,smooth,variable=\t]
      plot ({0.5*(cos(\t)-3*sin(\t))},{0.5*sin(\t)});
    \draw[thick,blue!55,domain=0:360,samples=90,smooth,variable=\t]
      plot ({1.0*(cos(\t)-3*sin(\t))},{1.0*sin(\t)});
    \fill (0,0) circle (1.5pt);
    \draw[blue] (-1.049,0.433) -- (-1.3,1.45) node[above,font=\scriptsize] {$f=0.5$};
    \draw[blue!55] (-2.098,0.866) -- (-2.85,1.45) node[above,font=\scriptsize] {$f=2$};
    \node[font=\small] at (0,-2.45) {$f_{20}=2(x+3y)^2+2y^2$: ellipses};
  \end{scope}
  \begin{scope}[xshift=8.6cm]
    \begin{scope}
      \clip (-3.7,-1.95) rectangle (3.7,1.95);
      \draw[dashed,gray,domain=-3:3] plot ({-0.655*\x},{\x});
      \draw[dashed,gray,domain=-0.7:0.7] plot ({-5.345*\x},{\x});
      \foreach \s in {1,-1} {
        \draw[thick,red!75!black,domain=-2.6:2.6,samples=90,smooth,variable=\t]
          plot ({(\s)*cosh(\t)-3*0.4264*sinh(\t)},{0.4264*sinh(\t)});
        \draw[thick,orange!85!black,domain=-2.6:2.6,samples=90,smooth,variable=\t]
          plot ({sinh(\t)-3*(\s)*0.4264*cosh(\t)},{(\s)*0.4264*cosh(\t)});
      }
    \end{scope}
    \draw[->,gray] (-3.6,0) -- (3.6,0) node[right] {$x$};
    \draw[->,gray] (0,-1.9) -- (0,1.9) node[above] {$y$};
    \fill (0,0) circle (1.5pt);
    \node[red!75!black,font=\scriptsize,fill=white,inner sep=1pt] at (1.55,0.35) {$f=2$};
    \node[orange!85!black,font=\scriptsize,fill=white,inner sep=1pt] at (-2.45,1.25) {$f=-2$};
    \node[font=\small] at (0,-2.45) {$f_{7}=2(x+3y)^2-11y^2$: hyperbolas};
  \end{scope}
\end{tikzpicture}
\caption{Level curves of two quadratic forms that differ only in the corner entry. Left: for the
positive definite $A_{20}$ every level set $f=k>0$ is an ellipse around the minimum at the origin.
Right: for the indefinite $A_7$ the level sets $f=\pm2$ are hyperbolas; along the dashed lines
$x+3y=\pm\sqrt{11/2}\,y$ the form is zero, and between them it changes sign: a saddle point.}
\label{fig:u31-levels}
\end{figure}

\subsection{Second derivatives decide a minimum}

In one-variable calculus, $f'(x_0)=0$ makes $x_0$ a critical point and $f''(x_0)>0$ makes it a
minimum: the slope is increasing as we pass through. For $f(x_1,\dots,x_n)$ the first
derivatives must all vanish, and the role of $f''$ is played by the \emph{matrix of second
derivatives} (the Hessian)
\[
  H=\begin{bmatrix}\dfrac{\partial^2f}{\partial x_i\,\partial x_j}\end{bmatrix},
  \qquad\text{for two variables}\quad
  H=\begin{bmatrix}f_{xx}&f_{xy}\\f_{yx}&f_{yy}\end{bmatrix}.
\]
It is symmetric because mixed partial derivatives can be taken in either order: $f_{xy}=f_{yx}$.

\begin{proposition}[Second-derivative test]\label{prop:u31-hessian}
Let $f$ have continuous second derivatives near $\vb{x}_0$, with all first derivatives zero at
$\vb{x}_0$. If $H(\vb{x}_0)$ is positive definite, $\vb{x}_0$ is a strict local minimum. If
$H(\vb{x}_0)$ has a negative eigenvalue, $\vb{x}_0$ is not a minimum.
\end{proposition}
\begin{proof}[Sketch]
Taylor's formula gives $f(\vb{x}_0+\vb{h})=f(\vb{x}_0)+\tfrac12\vb{h}\T H(\vb{x}_0)\vb{h}+r(\vb{h})$ with
$r(\vb{h})/\|\vb{h}\|^2\to0$ as $\vb{h}\to\vb{0}$; the linear term is absent because the gradient is zero.
If $H$ is positive definite with smallest eigenvalue $\lambda_{\min}>0$, then
$\tfrac12\vb{h}\T H\vb{h}\ge\tfrac12\lambda_{\min}\|\vb{h}\|^2$ (expand $\vb{h}$ in the orthonormal
eigenvectors), which beats $r(\vb{h})$ for small $\vb{h}\ne\vb{0}$, so $f(\vb{x}_0+\vb{h})>f(\vb{x}_0)$. If
$H\vb{q}=\lambda\vb{q}$ with $\lambda<0$, moving along $\vb{h}=t\vb{q}$ gives
$f(\vb{x}_0+t\vb{q})-f(\vb{x}_0)=\tfrac12\lambda t^2\|\vb{q}\|^2+r(t\vb{q})<0$ for small $t\ne0$.
The full $\varepsilon$--$\delta$ details belong to multivariable calculus.
\end{proof}

For two variables the determinant test for $H$ is the familiar calculus rule: a minimum when
$f_{xx}>0$ and $f_{xx}f_{yy}>f_{xy}^2$. Linear algebra extends it to any number of variables,
where it is applied through pivots, determinants or eigenvalues of $H$. For a pure quadratic
$f(\vb{x})=\vb{x}\T A\vb{x}$ the Hessian is the constant matrix $2A$, so the local picture at the origin
is the global one: a bowl when $A$ is positive definite, a saddle when it is indefinite.

\begin{example}[Classifying critical points]\label{ex:u31-critical}
Let $f(x,y)=x^3+y^2-3xy$. Setting the first derivatives to zero,
\[
  f_x=3x^2-3y=0,\qquad f_y=2y-3x=0 .
\]
The second equation gives $y=\tfrac32x$; then $3x^2-\tfrac92x=0$, so $x=0$ or $x=\tfrac32$. The
critical points are $(0,0)$ and $(\tfrac32,\tfrac94)$. The Hessian is
$H=\left[\begin{smallmatrix}6x&-3\\-3&2\end{smallmatrix}\right]$.
\begin{itemize}
  \item At $(0,0)$: $H=\left[\begin{smallmatrix}0&-3\\-3&2\end{smallmatrix}\right]$ has
        $\det H=-9<0$, so one eigenvalue is negative: a saddle point, not a minimum.
  \item At $(\tfrac32,\tfrac94)$: $H=\left[\begin{smallmatrix}9&-3\\-3&2\end{smallmatrix}\right]$ has
        $D_1=9>0$ and $D_2=18-9=9>0$: positive definite, so a local minimum, with value
        $f=\tfrac{27}{8}+\tfrac{81}{16}-\tfrac{81}{8}=-\tfrac{27}{16}$.
\end{itemize}
\end{example}

\subsection{Ellipses, ellipsoids and principal axes}

Cut the bowl $z=\vb{x}\T A\vb{x}$ at height $1$. The cross-section $\vb{x}\T A\vb{x}=1$ is an ellipse in the plane,
an ellipsoid (a lopsided rugby ball) in $\R^3$. Cutting a saddle instead gives a hyperbola. The
spectral theorem says exactly how the ellipse sits.

\begin{theorem}[Principal axes]\label{thm:u31-axes}
Let $A=Q\Lambda Q\T$ be positive definite with orthonormal eigenvectors $\vb{q}_1,\dots,\vb{q}_n$. The
set $\vb{x}\T A\vb{x}=1$ is an ellipsoid whose axes point along the eigenvectors $\vb{q}_i$, with
half-lengths $1/\sqrt{\lambda_i}$.
\end{theorem}
\begin{proof}
In the coordinates $\vb{y}=Q\T\vb{x}$, which measure $\vb{x}$ along the perpendicular directions
$\vb{q}_i$ (a rotation, possibly with a reflection, which changes no shapes), the equation becomes
$\lambda_1y_1^2+\cdots+\lambda_ny_n^2=1$, the standard ellipsoid
$\sum y_i^2/(1/\sqrt{\lambda_i})^2=1$. Its axes are the $y_i$-axes, i.e.\ the $\vb{q}_i$ directions,
and the axis along $\vb{q}_i$ reaches $y_i=\pm1/\sqrt{\lambda_i}$.
\end{proof}

Large eigenvalues give short axes: the bowl rises steeply in that direction. Equal eigenvalues
give a round cross-section; two equal out of three give a rugby ball; the identity gives a
sphere.

\begin{example}[An ellipse tilted at \texorpdfstring{$45^\circ$}{45 degrees}]\label{ex:u31-ellipse}
For $A=\begin{bmatrix}5&4\\4&5\end{bmatrix}$ the equation $\vb{x}\T A\vb{x}=1$ reads $5x^2+8xy+5y^2=1$.
The eigenvalues are $9$ and $1$ (trace $10$, determinant $9$) with eigenvectors
$\vb{q}_1=(1,1)/\sqrt2$ and $\vb{q}_2=(1,-1)/\sqrt2$. So
\[
  5x^2+8xy+5y^2=9\Bigl(\frac{x+y}{\sqrt2}\Bigr)^2+1\Bigl(\frac{x-y}{\sqrt2}\Bigr)^2 ,
\]
and the ellipse has half-axis $1/\sqrt9=\tfrac13$ along $(1,1)$ and $1/\sqrt1=1$ along $(1,-1)$
(Figure~\ref{fig:u31-axes}). Check the endpoints: $(x,y)=\tfrac{1}{3\sqrt2}(1,1)$ gives
$\tfrac{5+8+5}{18}=1$, and $(x,y)=\tfrac{1}{\sqrt2}(1,-1)$ gives $\tfrac{5-8+5}{2}=1$.

For $K$ of Example~\ref{ex:u31-K}, $\vb{x}\T K\vb{x}=1$ is an ellipsoid with three different axes:
the longest, $1/\sqrt{2-\sqrt2}\approx1.31$, along $(1,\sqrt2,1)$; the middle one, $1/\sqrt2\approx0.71$,
along $(1,0,-1)$; the shortest, $1/\sqrt{2+\sqrt2}\approx0.54$, along $(1,-\sqrt2,1)$.
\end{example}

\begin{figure}[ht]\centering
\begin{tikzpicture}[scale=2.1,>=Stealth]
  \draw[->,gray] (-1.05,0) -- (1.1,0) node[right] {$x$};
  \draw[->,gray] (0,-0.95) -- (0,0.95) node[above] {$y$};
  \draw[thick,blue,domain=0:360,samples=100,smooth,variable=\t]
    plot ({0.70711*(cos(\t)/3+sin(\t))},{0.70711*(cos(\t)/3-sin(\t))});
  \draw[->,very thick,red!75!black] (0,0) -- (0.2357,0.2357);
  \node[red!75!black,right,font=\small] at (0.27,0.33) {$\tfrac13\,\vb{q}_1$, $\lambda_1=9$};
  \draw[->,very thick,orange!85!black] (0,0) -- (0.70711,-0.70711);
  \node[orange!85!black,right,font=\small] at (0.72,-0.66) {$1\cdot\vb{q}_2$, $\lambda_2=1$};
  \draw (0.06,0.06) -- (0.12,0) -- (0.06,-0.06);
  \node[blue,font=\small] at (-0.85,-0.62) {$5x^2+8xy+5y^2=1$};
\end{tikzpicture}
\caption{The principal axes of $\vb{x}\T A\vb{x}=1$ for $A=\left[\begin{smallmatrix}5&4\\4&5\end{smallmatrix}\right]$
point along the eigenvectors; the half-lengths are $1/\sqrt{\lambda}$, so the large eigenvalue $9$
gives the short axis $\tfrac13$.}
\label{fig:u31-axes}
\end{figure}

\subsection{Exercises}

\begin{exercise}[Products of positive definite matrices]\label{exr:u31-AB}
When $A$ and $B$ are symmetric positive definite, $AB$ need not be symmetric, but its
eigenvalues are still positive. Prove this: start from $AB\vb{x}=\lambda\vb{x}$ and take the dot product
with $B\vb{x}$.
\end{exercise}
\begin{solution}
From $AB\vb{x}=\lambda\vb{x}$ with $\vb{x}\ne\vb{0}$, take the dot product with $B\vb{x}$:
\[
  (B\vb{x})\T AB\vb{x}=\lambda\,(B\vb{x})\T\vb{x}=\lambda\,\vb{x}\T B\vb{x},
\]
using $B\T=B$. The left side is $\vb{y}\T A\vb{y}$ with $\vb{y}=B\vb{x}\ne\vb{0}$ ($B$ is invertible), which is
positive because $A$ is positive definite. On the right, $\vb{x}\T B\vb{x}>0$ because $B$ is positive
definite. Hence $\lambda=\vb{y}\T A\vb{y}/\vb{x}\T B\vb{x}>0$. (Such a $\lambda$ is automatically real: the
quotient is real.) For instance $A=\left[\begin{smallmatrix}2&1\\1&1\end{smallmatrix}\right]$,
$B=\left[\begin{smallmatrix}1&1/2\\1/2&1\end{smallmatrix}\right]$ give the nonsymmetric
$AB=\left[\begin{smallmatrix}2.5&2\\1.5&1.5\end{smallmatrix}\right]$, whose eigenvalues
$2\pm\frac{\sqrt{13}}{2}\approx3.80,\,0.20$ are positive.
\end{solution}

\begin{exercise}[The form of a nonsymmetric matrix]\label{exr:u31-form}
Find the quadratic form $f(x,y)$ of $A=\left[\begin{smallmatrix}1&5\\7&9\end{smallmatrix}\right]$. Is it
always positive, always negative, or neither?
\end{exercise}
\begin{solution}
$f(x,y)=\begin{bmatrix}x&y\end{bmatrix}\begin{bmatrix}x+5y\\7x+9y\end{bmatrix}=x^2+5xy+7xy+9y^2=x^2+12xy+9y^2$.
Only the sum $5+7$ of the off-diagonal entries matters, so this is also the form of the
symmetric matrix $\left[\begin{smallmatrix}1&6\\6&9\end{smallmatrix}\right]$, whose determinant
$9-36=-27$ is negative. Hence $f$ takes both signs: $f(1,0)=1>0$ and $f(2,-2)=4-48+36=-8<0$.
Completing the square, $f=(x+6y)^2-27y^2$.
\end{solution}

\begin{exercise}[A parameter in the corner]\label{exr:u31-param}
For which $c$ is $B=\left[\begin{smallmatrix}2&-1&-1\\-1&2&-1\\-1&-1&2+c\end{smallmatrix}\right]$ positive
definite? Positive semidefinite? Confirm with pivots and by completing squares, and find
$\vb{x}\ne\vb{0}$ with $\vb{x}\T B\vb{x}=0$ when $c=0$.
\end{exercise}
\begin{solution}
\emph{Determinants.} $D_1=2$, $D_2=4-1=3$, and expanding along row~1,
$D_3=2\bigl(2(2+c)-1\bigr)+1\bigl(-(2+c)-1\bigr)-1\bigl(1+2\bigr)=(6+4c)-(3+c)-3=3c$. So $B$ is
positive definite for $c>0$ and positive semidefinite for $c\ge0$ (for $c<0$, $\det B<0$ forces a
negative eigenvalue).

\emph{Pivots.} Add $\tfrac12$ row~1 to rows 2 and 3: rows become $(0,\tfrac32,-\tfrac32)$ and
$(0,-\tfrac32,\tfrac32+c)$; add the new row~2 to row~3: $(0,0,c)$. Pivots $2,\ \tfrac32,\ c$, which are
the ratios $2,\ 3/2,\ 3c/3$ of the determinants.

\emph{Squares.} $\vb{x}\T B\vb{x}=2x^2+2y^2+(2+c)z^2-2xy-2xz-2yz
=2\bigl(x-\tfrac12y-\tfrac12z\bigr)^2+\tfrac32(y-z)^2+cz^2$. For $c=0$ this vanishes when $y=z$ and
$x=\tfrac12(y+z)$, for example at $\vb{x}=(1,1,1)$; indeed every row of $B$ then sums to zero, so
$(1,1,1)$ is in the nullspace.
\end{solution}

\begin{exercise}[How big can the cross term be?]\label{exr:u31-b}
For which $b$ is $\left[\begin{smallmatrix}1&b\\b&9\end{smallmatrix}\right]$ positive definite? For
$b=4$, find a vector that makes the form negative.
\end{exercise}
\begin{solution}
Need $1>0$ and $9-b^2>0$, so $-3<b<3$. At $b=\pm3$ the matrix is only semidefinite:
$x^2\pm6xy+9y^2=(x\pm3y)^2$. For $b=4$, $x^2+8xy+9y^2=(x+4y)^2-7y^2$; choosing $x+4y=0$, $y=1$,
i.e.\ $(x,y)=(-4,1)$, gives $16-32+9=-7<0$.
\end{solution}
